\section{Strict depth, signed oscillation and minimal positive compensation}
\label{ado:sec:scope}
The real-order depth-two factorization suggests a structural question at
higher strict depth: can the initial zero moments account for every sign
change of a representing density? The continuation~\cite{ado:report}
answers this when the leading indices are positive integers and the last
index is positive real. Its construction gives a minimal polynomial
compensator, a complete upper-arc count and exact integral identities.
The ordinary positive transform supplies the geometry only after the
signed-density zeros and initial moment cancellation have been established.
The proof chain follows those dependencies here.
\newcommand{\adoNorm}[1]{\lVert#1\rVert}
\newcommand{\adoArg}{\operatorname{Arg}}
\subsection{Strict single-variable notation}
For an index of depth $d$, define
\begin{equation}\label{ado:eq:series}
 L_{\boldsymbol{s}}(z)=\Li_{s_1,\ldots,s_d}(z,1,\ldots,1)
 =\sum_{n_1>\cdots>n_d\ge1}
       \frac{z^{n_1}}{n_1^{s_1}\cdots n_d^{s_d}}
 =\sum_{n\ge1}c_{\boldsymbol{s}}(n)z^n,\qquad |z|<1.
\end{equation}
Our parameter domain is
\begin{equation}\label{ado:eq:domain}
 s_1,\ldots,s_{d-1}\in\N_{\ge1},\qquad s_d=b>0.
\end{equation}
At depth one, this means any $b>0$. Indices are strictly decreasing, not
weakly decreasing. Only the outer summation variable receives $z^{n_1}$.
This convention agrees with the single-variable convention in the
manuscript and in the multiple-polylogarithm literature. The principal branch conventions are those fixed at the start of the book.

Write
\begin{equation}\label{ado:eq:leading}
 C_{\boldsymbol{s}}=c_{\boldsymbol{s}}(d)
 =\frac{1}{d^{s_1}(d-1)^{s_2}\cdots 2^{s_{d-1}}}>0.
\end{equation}
The empty product at depth one is one. The value of $C_{\boldsymbol{s}}$ is
independent of $b$, because the innermost index of the first admissible
chain is one. Also $c_{\boldsymbol{s}}(1)=\cdots=c_{\boldsymbol{s}}(d-1)=0$.
Throughout, principal continuation means continuation to
\[
 \Omega=\C\setminus[1,\infty).
\]
The word ``zero'' refers to a zero of a holomorphic function, unless the
imaginary part or a signed density is explicitly specified.

\subsection{Main result}
\begin{theorem}[All-depth compensation and zero geometry]
\label{ado:thm:main}
Let \eqref{ado:eq:domain} hold. There is a smooth real density
$\kappa_{\boldsymbol{s}}\in L^1(0,1)$ such that
\[
 \int_0^1u^{n-1}\kappa_{\boldsymbol{s}}(u)\dd u=c_{\boldsymbol{s}}(n),\qquad
 \adoNorm{\kappa_{\boldsymbol{s}}}_1\le2^{d-1}.
\]
It has exactly $d-1$ simple zeros $0<r_1<\cdots<r_{d-1}<1$, with
alternating signs and positive sign near one. Set
\[
 P_{\boldsymbol{s}}(u)=\prod_{j=1}^{d-1}(u-r_j),\qquad
 D_{\boldsymbol{s}}(z)=\prod_{j=1}^{d-1}(1-r_jz).
\]
Then $\dd\mu_{\boldsymbol{s}}=P_{\boldsymbol{s}}\kappa_{\boldsymbol{s}}\dd u$ is a positive
measure of mass $C_{\boldsymbol{s}}$, and
\begin{equation}\label{ado:eq:mainfactor}
 L_{\boldsymbol{s}}(z)=\frac{z^d}{D_{\boldsymbol{s}}(z)}H_{\boldsymbol{s}}(z),\qquad
 H_{\boldsymbol{s}}(z)=\int_0^1\frac{\dd\mu_{\boldsymbol{s}}(u)}{1-zu}.
\end{equation}
The polynomial $D_{\boldsymbol{s}}$ is the unique normalized polynomial of
smallest degree making $D_{\boldsymbol{s}}L_{\boldsymbol{s}}/z^d$ a positive
Hausdorff--Stieltjes transform. In particular, $L_{\boldsymbol{s}}$ has no zero
in $\Omega$ other than its zero of multiplicity $d$ at zero.

For every $0<\rho\le1$, there are exactly $d-1$ simple zeros
\[
 0<\theta_1(\rho)<\cdots<\theta_{d-1}(\rho)<\pi
\]
of $\theta\mapsto\Imag L_{\boldsymbol{s}}(\rho e^{\iu\theta})$.
With $\alpha_k=k\pi/d$ they satisfy
\begin{equation}\label{ado:eq:mainangles}
 \alpha_k-\arcsin(\rho\sin\alpha_k)
 <\theta_k(\rho)<\alpha_k.
\end{equation}
Every branch satisfies $\theta_k'(\rho)<0$ for $0<\rho\le1$ (with the
one-sided derivative at one). The imaginary part starts positive and
alternates sign at these zeros;
the value of $L_{\boldsymbol{s}}$ at the $k$th crossing is real with sign $(-1)^k$.
\end{theorem}
The proof is divided into operator, oscillation, compensation, and
angular-count steps below. None of those steps is inferred from the
finite verification suite.


\section{Two bounded operators on signed moment densities}
\label{ado:sec:operators}
For $f\in L^1(0,1)$ define its moments by
\[
 m_n(f)=\int_0^1u^{n-1}f(u)\dd u,\qquad n\ge1.
\]
We use the two operators
\begin{align}
 (\mathsf{A} f)(u)&=\int_u^1\frac{f(v)}v\dd v,
 \label{ado:eq:A}\\
 (\mathsf{B} f)(u)&=\int_0^u\frac{f(v)}{1-v}\dd v
              -\int_u^1\frac{f(v)}v\dd v.
 \label{ado:eq:B}
\end{align}
For every fixed $0<u<1$, both integrals in \eqref{ado:eq:B} are
absolutely convergent. In particular, this definition never subtracts
two divergent endpoint integrals. The operator $\mathsf{A}$ is the
outer-index-raising operator already present at depth two in the
manuscript~\cite{ado:report}. The companion $\mathsf{B}$ encodes a strict
prefix sum of moments.

\begin{lemma}[Norm, derivative, and moment identities]
\label{ado:lem:operators}
Both operators map $L^1(0,1)$ to itself and satisfy
\begin{equation}\label{ado:eq:opnorm}
 \adoNorm{\mathsf{A} f}_1\le\adoNorm f_1,\qquad
 \adoNorm{\mathsf{B} f}_1\le2\adoNorm f_1.
\end{equation}
Their moments are
\begin{equation}\label{ado:eq:opmoments}
 m_n(\mathsf{A} f)=\frac{m_n(f)}n,\qquad
 m_n(\mathsf{B} f)=\frac1n\sum_{j=1}^{n-1}m_j(f).
\end{equation}
If $f$ is continuous in the open interval, then
\begin{equation}\label{ado:eq:opderivatives}
 (\mathsf{A} f)'(u)=-\frac{f(u)}u,\qquad
 (\mathsf{B} f)'(u)=\frac{f(u)}{u(1-u)}.
\end{equation}
Moreover $\mathsf{A} f$ extends continuously to one with value zero, and
$m_1(\mathsf{B} f)=0$.
\end{lemma}
\begin{proof}
Tonelli's theorem applied to absolute values gives
\[
 \int_0^1\int_u^1\frac{|f(v)|}{v}\dd v\dd u
 =\int_0^1|f(v)|\dd v.
\]
Similarly,
\[
 \int_0^1\int_0^u\frac{|f(v)|}{1-v}\dd v\dd u
 =\int_0^1|f(v)|\dd v.
\]
These identities prove the two operator bounds and justify every
following application of Fubini to a signed integrand. For the first
operator,
\[
 m_n(\mathsf{A} f)=\int_0^1\frac{f(v)}v
       \left(\int_0^v u^{n-1}\dd u\right)\dd v
 =\frac{m_n(f)}n.
\]
The moment of the first term of $\mathsf{B} f$ is
\[
 \frac1n\int_0^1\frac{1-v^n}{1-v}f(v)\dd v
 =\frac1n\sum_{j=1}^n m_j(f).
\]
Subtracting $m_n(f)/n$ yields the strict prefix in
\eqref{ado:eq:opmoments}. The derivative formulas follow locally from
the fundamental theorem of calculus. Finally,
$\int_u^1f(v)/v\dd v\to0$ as $u\uparrow1$, because $1/v$ is bounded
near one and $f$ is integrable there.
\end{proof}

\subsection{Construction at arbitrary depth}
The positive starting density is
\begin{equation}\label{ado:eq:seed}
 f_b(u)=\frac{(-\log u)^{b-1}}{\Gamma(b)},\qquad b>0.
\end{equation}
Its mass is one and the Gamma integral gives $m_n(f_b)=n^{-b}$.
Define, with composition read from right to left,
\begin{equation}\label{ado:eq:kappaconstruction}
 \kappa_{\boldsymbol{s}}
 =\mathsf{A}^{s_1-1}\mathsf{B}\mathsf{A}^{s_2-1}\mathsf{B}\cdots
       \mathsf{A}^{s_{d-1}-1}\mathsf{B} f_b.
\end{equation}
There are $d-1$ occurrences of $\mathsf{B}$; all powers of $\mathsf{A}$ are
nonnegative integer powers. At depth one the right side is simply $f_b$.

\begin{proposition}[Moment realization and continuation]
\label{ado:prop:realization}
The density \eqref{ado:eq:kappaconstruction} is smooth on $(0,1)$,
belongs to $L^1(0,1)$, and satisfies
\begin{equation}\label{ado:eq:moments}
 m_n(\kappa_{\boldsymbol{s}})=c_{\boldsymbol{s}}(n),\qquad
 \adoNorm{\kappa_{\boldsymbol{s}}}_1\le2^{d-1}.
\end{equation}
The principal continuation is
\begin{equation}\label{ado:eq:signedtransform}
 L_{\boldsymbol{s}}(z)=z\int_0^1\frac{\kappa_{\boldsymbol{s}}(u)}{1-zu}\dd u,
 \qquad z\in\Omega.
\end{equation}
The density is the unique finite signed measure realizing these moments.
\end{proposition}
\begin{proof}
The coefficient recurrence for a nonempty tail $\boldsymbol{t}$ is
\[
 c_{(a,\boldsymbol{t})}(n)
 =\frac1{n^a}\sum_{j=1}^{n-1}c_{\boldsymbol{t}}(j).
\]
It is exactly the effect of $\mathsf{A}^{a-1}\mathsf{B}$ in
\eqref{ado:eq:opmoments}. Induction proves the moments and the norm
bound. The derivative identities prove local smoothness, starting from
\eqref{ado:eq:seed}. Expanding the denominator in the open disk is
justified by the integrable bound $|\kappa|/(1-|z|)$, and gives the
series \eqref{ado:eq:series}.

On a compact subset of $\Omega$, $|1-zu|$ is bounded away from zero
uniformly in $u\in[0,1]$. The signed integral is consequently holomorphic
there and continues the same germ. If two finite signed measures have
these moments, their difference integrates every polynomial to zero.
Uniform polynomial approximation on $[0,1]$, followed by uniqueness of
a finite measure from its continuous-function integrals, proves equality.
\end{proof}

\begin{corollary}[Ordinary convergence on the nontrivial unit circle]
\label{ado:cor:boundaryseries}
For \eqref{ado:eq:domain}, $c_{\boldsymbol{s}}(n)\to0$ and
\begin{equation}\label{ado:eq:coefficientvariation}
 \sum_{n\ge1}|c_{\boldsymbol{s}}(n+1)-c_{\boldsymbol{s}}(n)|
 \le\adoNorm{\kappa_{\boldsymbol{s}}}_1\le2^{d-1}.
\end{equation}
Thus the original series \eqref{ado:eq:series} converges at every
$|z|=1$, $z\ne1$, and agrees there with \eqref{ado:eq:signedtransform}.
\end{corollary}
\begin{proof}
Dominated convergence proves the coefficient limit. Also
\[
 |c(n+1)-c(n)|\le\int_0^1u^{n-1}(1-u)|\kappa(u)|\dd u.
\]
Summing the positive majorants gives \eqref{ado:eq:coefficientvariation}.
Summation by parts against the bounded partial sums of $z^n$ then proves
convergence when $|z|=1$, $z\ne1$. Abel's limiting theorem identifies its
sum with the radial limit of the disk series, which in turn equals the
continuous signed integral at that point.
\end{proof}
This conclusion concerns the parameter domain of this article. It is not
a claim that all of the manuscript's more general fractional-order
boundary series converge ordinarily.

\section{Exact oscillation and two interlacing laws}
\label{ado:sec:oscillation}
The central mechanism is a balance between moment orthogonality and
Rolle's theorem. Vanishing moments force sign changes; the two derivative
identities prevent extra zeros. This proof does not assume a
variation-diminishing theorem for fractional integration.

\begin{lemma}[Vanishing moments force sign changes]
\label{ado:lem:momentchanges}
Suppose a nonzero continuous real $f\in L^1(0,1)$ has finitely many
zeros and
\[
 \int_0^1u^j f(u)\dd u=0\qquad(0\le j\le m-1).
\]
Then $f$ has at least $m$ sign changes in $(0,1)$.
\end{lemma}
\begin{proof}
If its sign changes were $t_1,\ldots,t_k$ with $k<m$, the polynomial
$Q(u)=\prod_{j=1}^k(u-t_j)$ would have degree at most $m-1$ and
$Qf$ would have one fixed nonzero sign except at finitely many points.
Its integral would be nonzero, contradicting orthogonality. A zero at
which no sign change occurs is not included in $Q$.
\end{proof}

\begin{theorem}[Exact density zeros]
\label{ado:thm:oscillation}
For \eqref{ado:eq:domain}, $\kappa_{\boldsymbol{s}}$ has exactly $d-1$
distinct simple zeros in $(0,1)$. Its sign alternates at those zeros,
and it is positive near one.
\end{theorem}
\begin{proof}
At depth one, $f_b$ is strictly positive and has no zero. Suppose first
that a density $f$ at depth $d-1$ has exactly $d-2$ simple zeros.
By \eqref{ado:eq:opderivatives}, the derivative of $\mathsf{B} f$ has
exactly those zeros. Rolle's theorem therefore bounds the number of
distinct zeros of $\mathsf{B} f$ by $d-1$. This also rules out an infinite
set of zeros: any $d$ distinct zeros would already force too many
zeros of the derivative.

The first $d-1$ moments of the new density vanish, by
\eqref{ado:eq:opmoments}. Lemma~\ref{ado:lem:momentchanges} supplies
at least $d-1$ sign changes, so the upper bound is attained. There is
one derivative zero between each pair of consecutive new zeros, using
all $d-2$ derivative zeros. None remains at a new zero. Hence every
new zero is simple.

Now suppose $f$ is a depth-$d$ density with exactly $d-1$ simple
zeros, and put $g=\mathsf{A} f$. We have $g(1)=0$. If $g$ had $k$
interior zeros, Rolle's theorem between successive interior zeros and
between the largest one and the endpoint one would give at least $k$
zeros of $g'=-f/u$. Thus $k\le d-1$. Its first $d-1$ moments still
vanish, so the same lower bound gives exactly $d-1$ sign changes.
The derivative zeros are now all allocated to those successive
intervals, including the last interval ending at one. Consequently the
interior zeros of $g$ are simple. Repeating this argument for each
occurrence of $\mathsf{A}$ and $\mathsf{B}$ proves the zero count.

Let $P(u)$ be the monic polynomial whose roots are the resulting
$d-1$ zeros. The product $P\kappa_{\boldsymbol{s}}$ has a fixed strict sign
off its roots. Orthogonality to degrees at most $d-2$ gives
\[
 \int_0^1P(u)\kappa_{\boldsymbol{s}}(u)\dd u
 =\int_0^1u^{d-1}\kappa_{\boldsymbol{s}}(u)\dd u=C_{\boldsymbol{s}}>0.
\]
Thus $P\kappa_{\boldsymbol{s}}$ is positive, and since $P$ is positive to
the right of its largest root, so is $\kappa_{\boldsymbol{s}}$.
\end{proof}

\begin{corollary}[Index raising and depth raising interlace]
\label{ado:cor:interlacing}
For $d\ge2$, if $f=\kappa_{\boldsymbol{s}}$ has roots $r_1<\cdots<r_{d-1}$ and
$\mathsf{A} f$ has roots $q_1<\cdots<q_{d-1}$, then
\begin{equation}\label{ado:eq:Ainterlace}
 q_1<r_1<q_2<r_2<\cdots<q_{d-1}<r_{d-1}<1.
\end{equation}
Thus raising the first index by one moves every density zero strictly
toward zero. For $d\ge3$, if a depth-$(d-1)$ density has roots
$t_1<\cdots<t_{d-2}$ and $\mathsf{B} f$ has roots $r_1<\cdots<r_{d-1}$,
then
\begin{equation}\label{ado:eq:Binterlace}
 r_1<t_1<r_2<\cdots<t_{d-2}<r_{d-1}.
\end{equation}
\end{corollary}
\begin{proof}
At $d=2$, depth raising creates one interior zero and there is no
previous zero to interlace. For the stated ranges, these are precisely the
derivative-zero allocations in the preceding
proof. In the $\mathsf{A}$ case the last interval is $(q_{d-1},1)$; in the
$\mathsf{B}$ case all intervals are between consecutive interior zeros.
The number of derivative zeros equals the number of allocated intervals,
so the inequalities are strict and exhaustive.
\end{proof}

\subsection{Positive seeds, terminal mixtures, and a Lerch extension}
The preceding proof used only positivity, smoothness, integrability, and
the positive mass of the initial density. The explicit form
\eqref{ado:eq:seed} was needed to identify its moments as $n^{-b}$.

\begin{theorem}[Positive-seed extension]
\label{ado:thm:seedextension}
Let $f$ be a smooth strictly positive density in $L^1(0,1)$, of mass
$M_0>0$. Replace $f_b$ by $f$ in \eqref{ado:eq:kappaconstruction}.
All density-zero, interlacing, compensation, slit-plane nonvanishing,
and angular-count conclusions of this article remain valid. The norm
bound becomes $2^{d-1}M_0$, and the mass of the positive compensated
measure is
\[
 \frac{M_0}{d^{s_1}(d-1)^{s_2}\cdots2^{s_{d-1}}}.
\]
\end{theorem}
\begin{proof}
The operator identities still hold. After $d-1$ strict prefix steps,
the first $d-1$ moments vanish and the first nonzero moment is the
stated positive number. The proof of Theorem~\ref{ado:thm:oscillation}
therefore applies unchanged. The arguments in the next two sections
use only these moment and oscillation properties, and so establish the
remaining assertions as well.
\end{proof}
In particular, for fixed leading indices, every finite positive linear
combination
\begin{equation}\label{ado:eq:terminalmixture}
 \sum_{j=1}^N a_j L_{s_1,\ldots,s_{d-1},b_j}(z),
 \qquad a_j>0,\quad b_j>0,
\end{equation}
has the same nonvanishing and exact angular-count properties. Its seed
is $\sum a_jf_{b_j}$. This is a special convex-closure statement; sums
of arbitrary zero-free holomorphic functions need not be zero-free.

For a second application, take $\lambda>-1$ and
\[
 f_{b,\lambda}(u)=u^\lambda
                \frac{(-\log u)^{b-1}}{\Gamma(b)}.
\]
Its moments are $(n+\lambda)^{-b}$ and its mass is
$(1+\lambda)^{-b}$. The resulting strict shifted-terminal family is
\begin{equation}\label{ado:eq:lerchfamily}
 L^{(\lambda)}_{s_1,\ldots,s_{d-1};b}(z)
 =\sum_{n_1>\cdots>n_d\ge1}
 \frac{z^{n_1}}{n_1^{s_1}\cdots n_{d-1}^{s_{d-1}}(n_d+\lambda)^b}.
\end{equation}
It has the same minimal compensator degree $d-1$ and exactly $d-1$
angular zeros for $0<\rho\le1$. At depth one it is the usual
$z$-multiple of a Lerch series. No assertion is made here about shifts
in all denominators simultaneously.

\section{Canonical positive compensation and slit-plane nonvanishing}
\label{ado:sec:compensation}
Fix $\boldsymbol{s}$ and abbreviate $\kappa=\kappa_{\boldsymbol{s}}$,
$P=P_{\boldsymbol{s}}$, $D=D_{\boldsymbol{s}}$, and $C=C_{\boldsymbol{s}}$. Empty products
at depth one are one. The preceding section proves that
\begin{equation}\label{ado:eq:positiveweight}
 P(u)\kappa(u)>0
 \quad\text{for all }u\in(0,1)\setminus\{r_1,\ldots,r_{d-1}\}.
\end{equation}
Since $P$ is bounded on $[0,1]$, the positive measure
$\dd\mu=P\kappa\dd u$ is finite.

\begin{lemma}[Polynomial cancellation]
\label{ado:lem:polycancel}
Suppose $m_1(\kappa)=\cdots=m_{d-1}(\kappa)=0$. For any real
polynomial $E$ of degree at most $d-1$, the removable quotient at zero
satisfies
\begin{equation}\label{ado:eq:generalcompensator}
 E(z)\frac{L(z)}{z^d}
 =\int_0^1\frac{u^{d-1}E(1/u)\kappa(u)}{1-zu}\dd u.
\end{equation}
Here $u^{d-1}E(1/u)$ is a polynomial, including at $u=0$.
\end{lemma}
\begin{proof}
First subtract the initial $d-1$ zero moments in
\eqref{ado:eq:signedtransform} to obtain
\[
 L(z)=z^d\int_0^1\frac{u^{d-1}\kappa(u)}{1-zu}\dd u.
\]
For $0\le j\le d-1$, multiplying this last transform by $z^j$
and moving the factor $z^j$ to the density changes it to
$u^{d-1-j}\kappa(u)$. The difference is a finite linear combination
of the moments of degrees $d-1-j,\ldots,d-2$, all zero. Applying
this monomial identity to every coefficient of $E$ proves
\eqref{ado:eq:generalcompensator} in the disk and hence throughout $\Omega$.
\end{proof}

\begin{theorem}[Positive factorization and nonvanishing]
\label{ado:thm:factorization}
The measure $\mu$ has mass $C$, and
\begin{equation}\label{ado:eq:factorization}
 D(z)L(z)=z^d\int_0^1\frac{\dd\mu(u)}{1-zu}.
\end{equation}
Its transform $H$ has strictly positive imaginary part when
$\Imag z>0$, strictly negative imaginary part when $\Imag z<0$, and
$H(x)>0$ for every real $x<1$. The only zero of $L$ on $\Omega$ is its
zero of multiplicity $d$ at zero.
\end{theorem}
\begin{proof}
We have $u^{d-1}D(1/u)=P(u)$, so the preceding lemma gives
\eqref{ado:eq:factorization}. Its value after division by $z^d$ at
zero is $\int P\kappa=C$, as already shown in the oscillation proof.
For nonreal $z$,
\[
 \Imag H(z)=(\Imag z)\int_0^1
              \frac{u\dd\mu(u)}{|1-zu|^2}.
\]
The integral is strictly positive because the measure has positive
density except at finitely many interior points. For real $x<1$,
the integrand defining $H(x)$ is positive. Thus $H$ has no zero on
$\Omega$. All zeros of $D$ are $1/r_j>1$, outside $\Omega$, while
$D(0)=1$ and $H(0)=C>0$. The factorization proves the claimed
multiplicity and the global nonvanishing assertion.
\end{proof}

\begin{theorem}[Minimality and uniqueness of the compensator]
\label{ado:thm:minimality}
Among real polynomials $E$ with $E(0)>0$ for which $E(z)L(z)/z^d$
is the transform of a finite positive measure on $[0,1]$, the smallest
possible degree is $d-1$. Every such polynomial of degree $d-1$ is
$E(0)D(z)$. In particular the normalization $E(0)=1$ singles out $D$.
\end{theorem}
\begin{proof}
Existence in degree $d-1$ follows from
Theorem~\ref{ado:thm:factorization}. Suppose that $\deg E\le d-1$.
Lemma~\ref{ado:lem:polycancel} already realizes the transform as the
finite signed measure with density
\[
 Q_E(u)\kappa(u),\qquad Q_E(u)=u^{d-1}E(1/u).
\]
Uniqueness of finite Hausdorff moments implies that any purported
positive representing measure must be this signed measure. Its density
must therefore be nonnegative almost everywhere. Since $\kappa$
changes sign at every $r_j$, continuity forces $Q_E(r_j)=0$.
As $r_j\ne0$, these are $d-1$ distinct roots $1/r_j$ of $E$.
Consequently $\deg E\ge d-1$. At equality those roots determine
$E$ up to its constant factor, which is fixed by $E(0)$.
\end{proof}
This is minimality of \emph{polynomial compensation for an ordinary
positive Stieltjes representation}. It is not minimal depth in an
algebra of numerical periods and not a minimal generalized Stieltjes
order. Those are different questions.

\subsection{Quantitative consequences}
For $x>0$, positivity and $0<u<1$ give
\begin{equation}\label{ado:eq:negativebounds}
 \frac{Cx^d}{(1+x)D(-x)}
 <(-1)^dL(-x)<\frac{Cx^d}{D(-x)}.
\end{equation}
Indeed $C/(1+x)<H(-x)<C$ and $D(-x)>0$. Without computing the
compensator roots one still has the coarser bounds
\begin{equation}\label{ado:eq:coarsebounds}
 C\left(\frac{x}{1+x}\right)^d
 <(-1)^dL(-x)<Cx^d,
\end{equation}
because $1\le D(-x)\le(1+x)^{d-1}$, with the relevant strictness
supplied by $H$ even at depth one.

Write $P(u)=\sum_{j=0}^{d-1}p_ju^j$ and
\begin{equation}\label{ado:eq:compensatedmoments}
 h_n=\int_0^1u^n\dd\mu(u)
     =\sum_{j=0}^{d-1}p_jc_{\boldsymbol{s}}(n+j+1),\qquad n\ge0.
\end{equation}
For all $n,k\ge0$,
\begin{equation}\label{ado:eq:completepositivity}
 (-1)^k\Delta^kh_n
 =\int_0^1u^n(1-u)^k\dd\mu(u)>0,
 \qquad \Delta h_n=h_{n+1}-h_n.
\end{equation}
Every finite Hankel matrix $(h_{i+j})_{0\le i,j\le m}$ is positive
definite. The same is true of $(h_{i+j+1})$ and
$(h_{i+j}-h_{i+j+1})$. To see this, pair each matrix with a nonzero
real vector and integrate the square of its polynomial against,
respectively, $\mu$, $u\mu$, or $(1-u)\mu$. Each integral is strictly
positive. In particular,
\[
 h_nh_{n+2}-h_{n+1}^2>0.
\]
These are exact nonlinear positivity constraints on the original nested
harmonic coefficients and the compensator roots. They should not be
mistaken for rational-coefficient identities among special values.

\section{Exactly \texorpdfstring{$d-1$}{d-1} angular zeros on every upper semicircle}
\label{ado:sec:angular}
The signed transform gives, for $0<\rho\le1$ and $0<\theta<\pi$,
\begin{equation}\label{ado:eq:angulartransform}
 \Imag L(\rho e^{\iu\theta})
 =\rho\sin\theta\,J_\rho(\cos\theta),\qquad
 J_\rho(c)=\int_0^1
       \frac{\kappa(u)}{1-2\rho cu+\rho^2u^2}\dd u.
\end{equation}
A bound based only on simple sign-change arguments would leave open the
possibility of tangencies. We therefore prove a multiplicity-sensitive
upper bound before proving existence.

\subsection{An elementary confluent Cauchy-kernel bound}
\begin{lemma}[Angular variation bound with multiplicities]
\label{ado:lem:variation}
Suppose $\kappa\in L^1(0,1)$ is continuous, has exactly $m$ sign
changes at $r_1<\cdots<r_m$, and has a nonzero fixed sign on each
complementary interval. Then $J_\rho$ has at most $m$ zeros in
$(-1,1)$, counted with analytic multiplicity, for every $0<\rho\le1$.
\end{lemma}
\begin{proof}
The function
\begin{equation}\label{ado:eq:cauchymap}
 t(u)=\frac{1+\rho^2u^2}{2\rho u}
\end{equation}
is strictly decreasing on $(0,1)$, since
$t'(u)=(\rho^2u^2-1)/(2\rho u^2)<0$. Its values are at least one,
and it tends to infinity as $u\downarrow0$. Thus
\[
 J_\rho(c)=\int_0^1\frac{\kappa(u)}{2\rho u}
                         \frac1{t(u)-c}\dd u.
\]
The factor $\kappa(u)/u$ need not be integrable by itself; this causes
no difficulty because the displayed Cauchy factor is $O(u)$ at zero.
All the integrals below are interpreted with their complete kernels.

Suppose there are distinct zeros $c_1,\ldots,c_q\in(-1,1)$ with
positive multiplicities $\ell_1,\ldots,\ell_q$ summing to $m+1$.
If more zeros are available, retain any subcollection of total
multiplicity $m+1$. Form the proper rational function
\begin{equation}\label{ado:eq:variationrational}
 R(t)=\frac{\prod_{j=1}^m(t-t(r_j))}
             {\prod_{a=1}^q(t-c_a)^{\ell_a}}.
\end{equation}
Its partial fractions are linear combinations of
$(t-c_a)^{-h}$ for $1\le h\le\ell_a$. Differentiation of $J_\rho$
under the integral is valid locally for $c\in(-1,1)$; explicitly,
\[
 J_\rho^{(h-1)}(c)
 =(h-1)!\int_0^1\frac{\kappa(u)}{2\rho u}
                         \frac1{(t(u)-c)^h}\dd u.
\]
The selected zero multiplicities therefore imply
\begin{equation}\label{ado:eq:contradictoryintegral}
 \int_0^1\frac{\kappa(u)}{2\rho u}R(t(u))\dd u=0.
\end{equation}
This integral is absolutely convergent: near zero $R(t)=O(1/t)=O(u)$,
and near one its denominators stay positive and bounded away from zero.

But each factor $t(u)-t(r_j)$ changes sign exactly where $\kappa$
changes sign. The denominator of $R(t(u))$ is strictly positive, because
$t(u)\ge1>c_a$. Hence the integrand in
\eqref{ado:eq:contradictoryintegral} has one fixed nonzero sign except
at the $r_j$. Its integral cannot vanish. This contradiction proves the
bound and also excludes an identically zero $J_\rho$.
\end{proof}
The proof is a finite partial-fraction argument, including repeated poles.
No unproved strict-total-positivity assertion is being used.

\subsection{Positive compensation supplies the missing zeros}
For $0<\theta<\pi$ and $0\le u\le1$, put
\[
 \beta_u(\theta)=-\adoArg(1-\rho u e^{\iu\theta}).
\]
For $0<u<1$ one has
\begin{equation}\label{ado:eq:betasector}
 0<\beta_u(\theta)<\beta_1(\theta)<\frac\pi2.
\end{equation}
The angle is strictly increasing in $u$, since differentiation with
respect to $u$ gives $\rho\sin\theta/|1-\rho u e^{\iu\theta}|^2>0$.
A positive mixture of the resolvents $(1-zu)^{-1}$ stays strictly inside
this sector. Consequently, with $H$ as in
\eqref{ado:eq:factorization},
\begin{equation}\label{ado:eq:Hsector}
 0<\adoArg H(\rho e^{\iu\theta})<\beta_1(\theta).
\end{equation}
For example, the second inequality follows by rotating each resolvent
through $-\beta_1$ and observing that its imaginary part is strictly
negative for $0<u<1$. All rays lie in an angle of width less than
$\pi/2$, so no argument ambiguity arises.

\begin{theorem}[Complete angular count and signs]
\label{ado:thm:angular}
For \eqref{ado:eq:domain} and every $0<\rho\le1$,
$\Imag L(\rho e^{\iu\theta})$ has exactly $d-1$ simple zeros in
$0<\theta<\pi$. They have the signs and bounds stated in
Theorem~\ref{ado:thm:main}.
\end{theorem}
\begin{proof}
A continuous unwrapped argument of $L$ along the upper arc is
\begin{equation}\label{ado:eq:unwrappedphase}
 \Phi(\theta)=d\theta+
       \sum_{j=1}^{d-1}\beta_{r_j}(\theta)
       +\adoArg H(\rho e^{\iu\theta}).
\end{equation}
The factorization proves this formula exactly; $\Phi$ is not reduced
modulo $2\pi$. The sector bounds imply
\begin{equation}\label{ado:eq:phasebounds}
 d\theta<\Phi(\theta)<d\bigl(\theta+\beta_1(\theta)\bigr)<d\pi.
\end{equation}
At the other end of the arc, $\Phi(\theta)\to d\pi$ as
$\theta\uparrow\pi$, since all resolvents approach positive real
values and $H(-\rho)>0$.

If $\rho<1$, then $\Phi(\theta)\to0$ as $\theta\downarrow0$.
The case $\rho=1$ requires more care: $L(1)$ may be singular, and
there is no reason to assign it argument zero. Instead, each fixed
$r_j<1$ gives $\beta_{r_j}(\theta)\to0$, while
\eqref{ado:eq:Hsector} gives
\[
 \limsup_{\theta\downarrow0}\adoArg H(e^{\iu\theta})\le\pi/2.
\]
It follows that $0<\Phi(\theta)<\pi$ for sufficiently small positive
$\theta$. This weaker endpoint statement is all that is needed.

By continuity, the phase crosses each of $\pi,2\pi,\ldots,(d-1)\pi$.
These yield at least $d-1$ distinct zeros of the imaginary part. On the
other hand, Lemma~\ref{ado:lem:variation} with $m=d-1$, together
with \eqref{ado:eq:angulartransform}, gives at most $d-1$ zeros
counted with multiplicities. The change of variable $c=\cos\theta$
has nonzero derivative in the open arc, and the prefactor
$\rho\sin\theta$ is positive. Thus multiplicities are preserved.
Every zero is simple, and there is exactly one crossing of each listed
phase level. The levels must occur in their stated order, since a
reversal would force a repeated crossing of an earlier level.

The initial phase lies in $(0,\pi)$, so the imaginary part starts
positive. Simplicity forces alternating signs thereafter. At the
$k$th zero, the phase is $k\pi$ and the modulus is nonzero, so the real
value has sign $(-1)^k$.

For the location bounds, write $\alpha_k=k\pi/d$. At a zero,
\eqref{ado:eq:phasebounds} gives
\[
 \theta_k<\alpha_k<\theta_k+\beta_1(\theta_k).
\]
The function $h(\theta)=\theta+\beta_1(\theta)$ is strictly
increasing, because
\[
 h'(\theta)=\frac{1-\rho\cos\theta}
                  {1-2\rho\cos\theta+\rho^2}>0.
\]
For $\rho<1$, solving $h(\theta)=\alpha$ gives
$\theta=\alpha-\arcsin(\rho\sin\alpha)$. At $\rho=1$ the same
expression is the solution when $\alpha>\pi/2$; when
$\alpha\le\pi/2$ it is zero, and the strict lower bound follows
from $\theta_k>0$. This proves \eqref{ado:eq:mainangles}.
\end{proof}
The argument proves the needed crossings without asserting that
$\Phi$ is globally monotone. Such an assertion is unnecessary here.

\begin{corollary}[Analytic motion and no interior collisions]
\label{ado:cor:analyticbranches}
For fixed integer leading indices, each angular branch is real analytic
in $(\rho,b)\in(0,1)\times(0,\infty)$, and extends real analytically
across $\rho=1$ locally at each of its nontrivial upper-arc zeros.
The density roots $r_j$ are real analytic in $b>0$. No two of the
angular branches collide for $0<\rho\le1$.
\end{corollary}
\begin{proof}
On compact subsets of $\Real b>0$, the seed $f_b$ is holomorphic as
an $L^1$-valued function; factors $|\log(-\log u)|^j$ arising from
parameter differentiation have an integrable common majorant on smaller
compact subsets. The bounded operators preserve this holomorphic
parameter dependence. The signed transform is locally holomorphic in
$z\in\Omega$ and in $b$. Simplicity of its angular zeros permits the
real analytic implicit-function theorem. The unit-radius zeros remain
away from the cut, so the same local argument crosses $\rho=1$.
The density derivative formulas give joint real analyticity in $u,b$
in the open interval, and their simple roots obey the same theorem.
Distinctness follows from Theorem~\ref{ado:thm:angular}.
\end{proof}
The local extension across radius one does not imply that the global
count persists at all larger radii. For $L_{1,1}=\log^2(1-z)/2$, an
upper-half-plane imaginary zero requires $|1-z|=1$, so its angle obeys
$\cos\theta=\rho/2$. There are no open-upper-arc zeros when $\rho>2$.

\subsection{Strict radial motion of every angular branch}
The compensation yields more than a zero count. It also turns a
monotonicity property of ordinary positive Stieltjes transforms into
strict motion of all the strict-depth angular zeros.

\begin{lemma}[Radial argument monotonicity for positive transforms]
\label{ado:lem:radialPick}
Let $\nu$ be a nonzero finite positive measure on $[0,1]$ and
$Q(z)=\int(1-uz)^{-1}\dd\nu(u)$. For $0<\theta<\pi$ and
$0<\rho\le1$,
\begin{equation}\label{ado:eq:radialPick}
 \frac{\partial}{\partial\rho}\adoArg Q(\rho e^{\iu\theta})\ge0.
\end{equation}
The argument is chosen continuously from $Q(0)>0$.
\end{lemma}
\begin{proof}
First take a finite measure with positive weights at distinct points
$0<u_1<\cdots<u_n<1$. The rational function
$\sum_j w_j/(t-u_j)$ is strictly decreasing between consecutive poles,
with limits $+\infty$ and $-\infty$ at the two ends. Its numerator
therefore has exactly one root $v_j\in(u_j,u_{j+1})$ in each such
interval. Factoring the numerator gives
\[
 Q(z)=Q(0)\frac{\prod_{j=1}^{n-1}(1-v_jz)}
                      {\prod_{j=1}^n(1-u_jz)}.
\]
Consequently the radial derivative of its argument is
\[
 \sin\theta\left\{f(u_1)+
          \sum_{j=1}^{n-1}\bigl(f(u_{j+1})-f(v_j)\bigr)\right\},
 \qquad f(t)=\frac{t}{1-2\rho t\cos\theta+\rho^2t^2}.
\]
Every summand is nonnegative, because
\[
 f'(t)=\frac{1-\rho^2t^2}
            {(1-2\rho t\cos\theta+\rho^2t^2)^2}\ge0
 \qquad(0\le t\le1).
\]

Approximate an arbitrary finite positive $\nu$ weakly by such finite
measures, with their masses converging to $\nu([0,1])$. For each
compact subset of $\Omega$, the integrands defining the transform and
its first complex derivative are bounded and continuous uniformly in
$u\in[0,1]$. The transforms and derivatives converge locally uniformly.
Also $Q(\rho e^{\iu\theta})\ne0$: it has positive imaginary part
unless the measure is supported at zero, in which case it is a positive
constant. Passing to the limit in
$\Imag(e^{\iu\theta}Q'(z)/Q(z))$ proves the nonnegative derivative.
The same argument works at $\rho=1$ because $\theta$ is strictly
between zero and $\pi$.
\end{proof}

\begin{theorem}[Strict inward angular motion]
\label{ado:thm:radialmotion}
For $d\ge2$ and \eqref{ado:eq:domain}, every angular branch satisfies
\begin{equation}\label{ado:eq:strictmotion}
 \frac{\dd\theta_k}{\dd\rho}<0,
 \qquad 0<\rho\le1,\quad 1\le k\le d-1.
\end{equation}
The derivative at one is the one-sided derivative, or equivalently the
derivative of the local continuation. The same result holds for the
positive-seed and shifted-terminal families of
Theorem~\ref{ado:thm:seedextension}.
\end{theorem}
\begin{proof}
At fixed $\theta$, differentiate the unwrapped phase
\eqref{ado:eq:unwrappedphase} with respect to the radius:
\[
 \Phi_\rho
 =\sum_{j=1}^{d-1}
      \frac{r_j\sin\theta}{1-2\rho r_j\cos\theta+\rho^2r_j^2}
   +\partial_\rho\adoArg H(\rho e^{\iu\theta})>0.
\]
The first sum is strictly positive since $d\ge2$ and $r_j>0$; the
last term is nonnegative by Lemma~\ref{ado:lem:radialPick}.
At the unique simple crossing of the phase level $k\pi$, one has
$\Phi_\theta>0$. Indeed the phase starts below that level, ends
above it, and crosses it exactly once; simplicity excludes a zero
phase derivative there. Implicit differentiation of
$\Phi(\theta_k(\rho),\rho)=k\pi$ now gives
\[
 \theta_k'(\rho)=-\frac{\Phi_\rho}{\Phi_\theta}<0.
\]
All ingredients also hold for a general positive seed.
\end{proof}
For depth two, angular decrease alone does \emph{not} settle the
higher-leading-order question about $\eta(\rho)=\cos\theta(\rho)/\rho$: the derivative of that
quotient has a second term, and its sign is not implied by
$\theta'(\rho)<0$. The two kinds of radial monotonicity must remain
separate in the integration ledger.

\section{Small-radius and large-outer-index expansions}
\label{ado:sec:asymptotics}
The global count removes an ambiguity from local expansions: the zeros
found near the first Fourier-mode zeros are the entire list of angular
zeros, not just selected branches.

\begin{theorem}[Two-term radial expansion]
\label{ado:thm:radialexpansion}
Fix an index satisfying \eqref{ado:eq:domain} with $d\ge2$, and put
\[
 p=\frac{c_{\boldsymbol{s}}(d+1)}{C_{\boldsymbol{s}}},\qquad
 q=\frac{c_{\boldsymbol{s}}(d+2)}{C_{\boldsymbol{s}}},\qquad
 \alpha_k=\frac{k\pi}{d}.
\]
The $k$th angular branch extends real analytically to $\rho=0$, and
\begin{align}
 \theta_k(\rho)={}&\alpha_k
 -\frac{p}{d}\sin\alpha_k\,\rho \notag\\
 &+\left(\frac{d+1}{d^2}p^2-\frac{2q}{d}\right)
       \sin\alpha_k\cos\alpha_k\,\rho^2
 +O_{\boldsymbol{s}}(\rho^3).
 \label{ado:eq:radialexpansion}
\end{align}
In particular, every branch initially moves strictly toward smaller
angles as the radius increases from zero.
\end{theorem}
\begin{proof}
Divide the imaginary part of the disk series by $C_{\boldsymbol{s}}\rho^d$.
The resulting analytic equation is
\[
 0=\sin(d\theta)+p\rho\sin((d+1)\theta)
                   +q\rho^2\sin((d+2)\theta)+O(\rho^3).
\]
At $(\rho,\theta)=(0,\alpha_k)$ its $\theta$ derivative is
$d(-1)^k\ne0$, so the analytic implicit-function theorem applies.
There are $d-1$ such branches; for small positive radius they exhaust
the global list from Theorem~\ref{ado:thm:angular}.

Write $\theta=\alpha_k+A\rho+B\rho^2+O(\rho^3)$ and divide by
$(-1)^k$. The coefficients of $\rho$ and $\rho^2$ are respectively
\[
 dA+p\sin\alpha_k,
 \qquad dB+(d+1)pA\cos\alpha_k+q\sin(2\alpha_k).
\]
Solving gives \eqref{ado:eq:radialexpansion}. The first derivative is
negative because $p>0$ and $0<\alpha_k<\pi$.
\end{proof}
Higher Taylor coefficients are obtained by the same finite coefficient
recursion. At each order the new unknown occurs with coefficient
$d(-1)^k$, while all remaining terms depend only on earlier
coefficients and finitely many moments. This is an exact algebraic
recurrence, not a proof of any global radial sign.

\subsection{Uniformity as the first index increases}
Fix $d\ge2$ and positive integer middle indices
$s_2,\ldots,s_{d-1}$. Let the first index be $a\in\N$ and the
last be $b>0$. Define
\[
 T_N(b)=\sum_{N\ge n_2>\cdots>n_d\ge1}
       \frac1{n_2^{s_2}\cdots n_{d-1}^{s_{d-1}}n_d^b},
 \qquad T_*=T_{d-1}(b)>0.
\]
The latter number is a single-chain value and is independent of $b$.
At depth two, $T_N(b)=H_N^{(b)}$ and $T_*=1$.

\begin{theorem}[First displacement, uniformly in the final order]
\label{ado:thm:largeouter}
Set $\lambda=d/(d+1)$ and $\mu=d/(d+2)$. As $a\to\infty$ through
positive integers, uniformly for $b>0$, $0\le\rho\le1$, and
$1\le k\le d-1$,
\begin{equation}\label{ado:eq:largeouter}
 \theta_k(\rho)
 =\alpha_k-\frac{T_d(b)}{dT_*}\sin\alpha_k\,
                      \rho\lambda^a
       +O\bigl(\rho^2\mu^a\bigr).
\end{equation}
At $\rho=0$, the branch is interpreted by its analytic extension.
The implied constant and the starting threshold for $a$ may depend on
the fixed depth and middle indices, but not on $b$, $\rho$, or $k$.
\end{theorem}
\begin{proof}
The normalized imaginary part is
\begin{equation}\label{ado:eq:outereq}
 \sin(d\theta)+\frac{T_d(b)}{T_*}\rho\lambda^a
                   \sin((d+1)\theta)+R(\theta),
\end{equation}
where
\[
 R(\theta)=\sum_{n\ge d+2}\frac{T_{n-1}(b)}{T_*}
        \rho^{n-d}\left(\frac dn\right)^a\sin(n\theta).
\]
All inner denominators are at least one, so
$0\le T_{n-1}(b)\le n^{d-1}$ uniformly for $b>0$. Choose the fixed
integer $a_0=d+3$. For $a\ge a_0$, termwise differentiation is
absolutely justified and
\begin{equation}\label{ado:eq:outertail}
 |R(\theta)|+|R'(\theta)|\le K\rho^2\mu^a,
 \quad
 K=\frac2{T_*}\sum_{n\ge d+2}n^d
                        \left(\frac{d+2}{n}\right)^{a_0}<\infty.
\end{equation}
Also $T_d(b)/T_*\le d^{d-1}/T_*$, independently of $b$.

On disjoint fixed small neighborhoods of the $\alpha_k$, the derivative
of the first term of \eqref{ado:eq:outereq} is bounded away from zero.
Its perturbation tends to zero uniformly by \eqref{ado:eq:outertail}.
Evaluating at $\alpha_k\pm K_1\rho\lambda^a$, for a sufficiently
large fixed $K_1$, gives opposite signs for all sufficiently large $a$.
This follows from the linear term $d(\theta-\alpha_k)$ and the
uniform bounds on the perturbations; $\mu/\lambda<1$. Thus each
neighborhood contains a unique root with displacement
$\delta=O(\rho\lambda^a)$. The global count identifies these roots
as all the branches in their established order.

Expand \eqref{ado:eq:outereq} at $\alpha_k$ and divide by $(-1)^k$.
The root equation becomes
\[
 0=d\delta+\frac{T_d(b)}{T_*}\rho\lambda^a\sin\alpha_k
     +O\bigl(\delta^2+\rho\lambda^a|\delta|
                         +\rho^2\mu^a\bigr).
\]
Because $\lambda^2<\mu$, the previously obtained displacement bound
makes the error $O(\rho^2\mu^a)$. Division by $d$ proves
\eqref{ado:eq:largeouter}. All constants used above are independent
of the final order.
\end{proof}
At depth two and unit radius this recovers the leading scale
$(2/3)^a$ already present in the manuscript's more detailed depth-two
expansion~\cite{ado:report}. The added content is the simultaneous
all-depth statement, the exact branch indexing, and uniformity over the
entire range $b>0$ and $0\le\rho\le1$.

\section{Explicit logistic kernels and positive integral identities}
\label{ado:sec:identities}
The compensator is generally implicit: its roots are the zeros of a
signed density. The all-ones family makes every part explicit and gives
an infinite collection of exact positive integral identities.

\subsection{The all-ones density and its roots}
Let $\boldsymbol{1}_d=(1,\ldots,1)$ and define
\[
 \ell(u)=\log\frac{u}{1-u},\qquad 0<u<1.
\]
The familiar identity
\begin{equation}\label{ado:eq:allonesLi}
 L_{\boldsymbol{1}_d}(z)=\frac{(-\log(1-z))^d}{d!}
\end{equation}
follows recursively from
$(1-z)L'_{\boldsymbol{1}_d}(z)=L_{\boldsymbol{1}_{d-1}}(z)$ and the vanishing initial
value at zero. We use it as a classical identity with a short proof,
not as a new reduction of multiple polylogarithms.

\begin{theorem}[Explicit signed kernels and logistic roots]
\label{ado:thm:logistickernels}
For every $d\ge1$,
\begin{align}
 \kappa_{\boldsymbol{1}_d}(u)
 &=\frac{\Imag(\ell(u)+\iu\pi)^d}{\pi d!}\notag\\
 &=\sum_{j=0}^{\lfloor(d-1)/2\rfloor}
 \frac{(-1)^j\pi^{2j}\ell(u)^{d-1-2j}}
 {(2j+1)!(d-1-2j)!}.
 \label{ado:eq:logistickernels}
\end{align}
Its increasing list of roots is
\begin{equation}\label{ado:eq:logisticroots}
 r_j=\frac1{1+\exp(\pi\cot(j\pi/d))},
 \qquad 1\le j\le d-1.
\end{equation}
In particular, $r_j+r_{d-j}=1$.
\end{theorem}
\begin{proof}
Consider the generating function
\begin{equation}\label{ado:eq:kernelgenerator}
 G(t,u)=e^{t\ell(u)}\frac{\sin\pi t}{\pi t}
       =\sum_{d\ge1}K_d(u)t^{d-1}.
\end{equation}
For $|\Real t|<1$, its integral at moment one is the beta integral
\[
 \int_0^1G(t,u)\dd u
 =\frac{\sin\pi t}{\pi t}\Gamma(1+t)\Gamma(1-t)=1.
\]
The reflection formula supplies the last equality, with removable values
understood at zero. Thus $\int K_1=1$ and $\int K_d=0$ for $d\ge2$.
Coefficient extraction is legitimate on a smaller complex neighborhood
of zero, since $u^{\Real t}(1-u)^{-\Real t}$ and its logarithmic
parameter derivatives have an integrable majorant.

Now $K_1=1$ and
\[
 \frac{\dd}{\dd u}G(t,u)=\frac{t}{u(1-u)}G(t,u),
\]
so $K_d'=K_{d-1}/(u(1-u))$. This is exactly the derivative of
$\mathsf{B} K_{d-1}$. Their difference is constant, and both have zero
integral for $d\ge2$, so $K_d=\mathsf{B} K_{d-1}$. Hence these are the
constructed all-ones kernels. Expanding the exponential and sine gives
\eqref{ado:eq:logistickernels}.

The imaginary part of $(\ell+\iu\pi)^d$ vanishes when
$\adoArg(\ell+\iu\pi)$ is one of $j\pi/d$. This gives
$\ell=\pi\cot(j\pi/d)$. Reversing the order of these real values
and solving $u/(1-u)=e^\ell$ yields \eqref{ado:eq:logisticroots}.
The symmetry follows from $\cot(\pi-x)=-\cot x$.
\end{proof}

\begin{proposition}[A continuous Mellin identity and Stirling moments]
\label{ado:prop:mellin}
For real $x>0$ and complex $t$ satisfying $-x<\Real t<1$,
\begin{equation}\label{ado:eq:mellin}
 \int_0^1u^{x-1}e^{t\ell(u)}\frac{\sin\pi t}{\pi t}\dd u
 =\frac{\Gamma(x+t)}{\Gamma(x+1)\Gamma(1+t)}.
\end{equation}
The values at removable singularities are obtained by continuation.
For every integer $n\ge1$,
\begin{equation}\label{ado:eq:stirlingmoments}
 \int_0^1u^{n-1}\kappa_{\boldsymbol{1}_d}(u)\dd u
 =\frac{\stirling{n}{d}}{n!},
\end{equation}
where $\stirling{n}{d}$ is the unsigned Stirling number of the first kind,
zero for $d>n$.
\end{proposition}
\begin{proof}
The integral equals
$\sin(\pi t)B(x+t,1-t)/(\pi t)$. The beta formula and Gamma
reflection simplify it to \eqref{ado:eq:mellin}. At integer $x=n$,
the right side is
\[
 \frac{(t+1)(t+2)\cdots(t+n-1)}{n!}.
\]
Its coefficient of $t^{d-1}$ is $\stirling{n}{d}/n!$, proving
\eqref{ado:eq:stirlingmoments} by \eqref{ado:eq:kernelgenerator}.
\end{proof}
These beta and Stirling identities provide an independent normalization
check on the general signed-moment construction.

\subsection{An explicit infinite family of compensated identities}
For the roots \eqref{ado:eq:logisticroots}, define
\[
 P_d(u)=\prod_{j=1}^{d-1}(u-r_j),\qquad
 D_d(z)=\prod_{j=1}^{d-1}(1-r_jz).
\]
The integrable function $P_d\kappa_{\boldsymbol{1}_d}$ is strictly positive
away from its finitely many zeros. Combining
\eqref{ado:eq:factorization} and \eqref{ado:eq:allonesLi} gives the
exact identity
\begin{equation}\label{ado:eq:generalidentity}
 \boxed{\quad
 \int_0^1\frac{P_d(u)\kappa_{\boldsymbol{1}_d}(u)}{1-zu}\dd u
 =\frac{D_d(z)}{d!}
       \left(\frac{-\log(1-z)}z\right)^d,
 \qquad z\in\Omega.
 \quad}
\end{equation}
At $z=0$ the removable value is $1/d!$. At $z=-1$ this becomes
\begin{equation}\label{ado:eq:negativeoneidentity}
 \int_0^1\frac{P_d(u)\kappa_{\boldsymbol{1}_d}(u)}{1+u}\dd u
 =\frac{D_d(-1)}{d!}\log^d2.
\end{equation}
The content is not a new formula for $L_{\boldsymbol{1}_d}$; it is the explicit
positive density, its canonical compensation, and the resulting
positive-integrand identities. They remain valid at every depth without
a numerical relation search.

\subsection{Depth three}
Put
\[
 h_3=\frac\pi{\sqrt3},\qquad
 q_3=\frac1{2+2\cosh h_3}.
\]
Then
\[
 P_3(u)=u^2-u+q_3,\qquad
 D_3(z)=1-z+q_3z^2,\qquad
 \kappa_{\boldsymbol{1}_3}(u)=\frac{\ell(u)^2-\pi^2/3}{2}.
\]
Both factors $P_3$ and $\ell^2-\pi^2/3$ are negative between the
same two roots and positive outside them. Hence their product is
nonnegative throughout the interval. Formula
\eqref{ado:eq:negativeoneidentity} gives
\begin{equation}\label{ado:eq:depththreeidentity}
 \boxed{
 \int_0^1\frac{(u^2-u+q_3)(\ell(u)^2-\pi^2/3)}{1+u}\dd u
 =\frac{2+q_3}{3}\log^3 2.}
\end{equation}
The companion normalization is
\begin{equation}\label{ado:eq:depththreenormalization}
 \int_0^1 (u^2-u+q_3)(\ell(u)^2-\pi^2/3)\dd u=\frac13.
\end{equation}
These identities are proved statements, not high-precision conjectures.

\subsection{Depth four}
Set
\[
 q_4=\frac1{2+2\cosh\pi},\qquad
 P_4(u)=\left(u-\frac12\right)(u^2-u+q_4).
\]
Now
\[
 D_4(z)=\left(1-\frac z2\right)(1-z+q_4z^2),\qquad
 \kappa_{\boldsymbol{1}_4}(u)=\frac{\ell(u)^3-\pi^2\ell(u)}6.
\]
The roots are the logistic images of $-\pi,0,\pi$, so again the
product $P_4\kappa_{\boldsymbol{1}_4}$ is nonnegative. The corresponding pair
of identities is
\begin{align}
 \int_0^1\frac{P_4(u)(\ell(u)^3-\pi^2\ell(u))}{1+u}\dd u
 &=\frac{3(2+q_4)}8\log^4 2,
 \label{ado:eq:depthfouridentity}\\
 \int_0^1P_4(u)(\ell(u)^3-\pi^2\ell(u))\dd u&=\frac14.
 \label{ado:eq:depthfournormalization}
\end{align}
All endpoint singularities are logarithmic powers and are integrable.
Thus no principal-value convention is implicit in these formulas.

\subsection{An explicit density-edge asymptotic}
The smallest all-ones density root gives a simple large-depth benchmark:
\begin{equation}\label{ado:eq:edgeasymptotic}
 \log r_1=-d+\frac{\pi^2}{3d}+\frac{\pi^4}{45d^3}
             +O(d^{-5})+O(e^{-d}).
\end{equation}
Indeed $\log r_1=-\pi\cot(\pi/d)
-\log(1+e^{-\pi\cot(\pi/d)})$, and the Taylor expansion of $\cot$
at zero gives the displayed terms. This concerns a density root, not
an angular zero. Symmetry supplies $1-r_{d-1}=r_1$.

\section{A divided-difference identity across depths and colors}
\label{logres:sec:main}
The explicit logistic density suggests an identity search that keeps the
depth parameter symbolic instead of recognizing one numerical value at a
time. Its beta moments already determine the one-resolvent transform.
Partial fractions then introduce several independent colors, while
coalescing those colors produces finite formulas at every resolvent order.
The resulting Gaussian logarithmic moments are exact polynomials in
$\pi$ and $\log2$, with no integer-relation premise.

For $-1<\Real t<1$ put
\[
 W_t(u)=\left(\frac{u}{1-u}\right)^t\frac{\sin\pi t}{\pi t},
 \qquad 0<u<1,
\]
using the removable value at $t=0$. This is the generating density
\eqref{ado:eq:kernelgenerator}. Its integral is one, even for complex $t$;
for real $-1<t<1$ it is positive. Denote divided differences by
$[z_1,\ldots,z_m]f$, with the continuous confluent interpretation when
nodes coincide.

\begin{theorem}[Several resolvents and their confluent limits]
\label{logres:thm:colors}
Let $z_1,\ldots,z_m\in\Omega$, where $\Omega=\C\setminus[1,\infty)$.
For distinct nodes and $m\ge2$,
\begin{equation}\label{logres:eq:colors}
 \int_0^1\frac{W_t(u)}{\prod_{j=1}^m(1-z_ju)}\,du
 =\frac1t[z_1,\ldots,z_m]
       \{z^{m-2}((1-z)^{-t}-1)\}.
\end{equation}
For $m=1$ the transform is
\begin{equation}\label{logres:eq:one}
 \int_0^1\frac{W_t(u)}{1-zu}\,du
       =\frac{(1-z)^{-t}-1}{tz}.
\end{equation}
All powers use the principal logarithm; the apparent singularities at
$t=0$, $z=0$, and coincident nodes are removable. In particular, for
$m\ge2$,
\begin{align}\label{logres:eq:confluent}
 \int_0^1\frac{W_t(u)}{(1-zu)^m}\,du
 &=(1-z)^{-t-m+1}\mathcal P_m(t,z),\notag\\
 \mathcal P_m(t,z)
 &=\sum_{j=0}^{m-2}\binom{m-2}{j}
       \frac{(t+1)_j}{(j+1)!}
       z^j(1-z)^{m-2-j},
\end{align}
where $(t+1)_j$ is the rising factorial and $(t+1)_0=1$.
\end{theorem}
\begin{proof}
The beta identity \eqref{ado:eq:mellin} gives
\[
 \int_0^1u^nW_t(u)\,du=\frac{(1+t)_n}{(n+1)!}.
\]
For $|z|<1$, expand the geometric kernel and sum the binomial series.
Its result is \eqref{logres:eq:one}. The integrable endpoint majorants
$u^{-1+\epsilon}(1-u)^{-1+\epsilon}$, on smaller parameter strips,
justify local holomorphy and continuation to the principal slit plane.
For distinct nonzero nodes, the exact partial fractions are
\[
 \frac1{\prod_j(1-z_ju)}
 =\sum_{j=1}^m
   \frac{z_j^{m-1}}{\prod_{k\ne j}(z_j-z_k)}\frac1{1-z_ju}.
\]
Integrating and using \eqref{logres:eq:one} gives the Lagrange formula
for \eqref{logres:eq:colors}. Continuity allows zero nodes.

When every node tends to $z$, the divided difference becomes the
$(m-1)$st derivative divided by $(m-1)!$. Apply the product rule to
$z^{m-2}((1-z)^{-t}-1)/t$. The term differentiating the second factor
zero times vanishes, because the polynomial has degree $m-2$.
For $j+1\ge1$ derivatives of the second factor one has
\[
 \frac{d^{j+1}}{dz^{j+1}}
       \frac{(1-z)^{-t}-1}{t}
       =(t+1)_j(1-z)^{-t-j-1}.
\]
Collecting the binomial factors proves \eqref{logres:eq:confluent}.
All confluent limits also follow directly from domination of the
original integral on compact subsets of $\Omega$.
\end{proof}

For example, the first three confluent cases are
\begin{align*}
 \mathcal P_2&=1,\\
 \mathcal P_3&=1+(t-1)z/2,\\
 \mathcal P_4&=(1-z)^2+(t+1)z(1-z)
                         +(t+1)(t+2)z^2/6.
\end{align*}
Thus the squared-resolvent identity has the particularly simple form
\[
 \int_0^1\frac{W_t(u)}{(1-zu)^2}\,du=(1-z)^{-t-1}.
\]
The several-color identity is symmetric in its nodes; repeated colors
are derivatives, rather than a separate guessed formula.

\begin{corollary}[Every Gaussian logarithmic moment]
\label{logres:cor:Gaussian}
For integers $m\ge2$ and $k\ge0$, define
\[
 I_{m,k}=\int_0^1\frac{\log^k(u/(1-u))}{(1-\iu u)^m}\,du,
 \qquad \lambda=\tfrac12\log2-\tfrac{\iu\pi}4.
\]
These are ordinary absolutely convergent integrals, and
\begin{equation}\label{logres:eq:Gaussian}
 I_{m,k}=\frac{k!}{(1-\iu)^{m-1}}
 [t^k]\left\{\mathcal P_m(t,\iu)
   \exp\left(-\lambda t+
       \sum_{j\ge1}\frac{\zeta(2j)}j t^{2j}\right)\right\}.
\end{equation}
In particular $I_{m,k}\in\Q[\iu,\pi,\log2]$ for every $m,k$.
For the squared resolvent,
\begin{align*}
 I_{2,0}&=\frac1{1-\iu},&
 I_{2,1}&=-\frac{\lambda}{1-\iu},\\
 I_{2,2}&=\frac{\lambda^2+\pi^2/3}{1-\iu},&
 I_{2,3}&=-\frac{\lambda^3+\pi^2\lambda}{1-\iu},\\
 I_{2,4}&=\frac{\lambda^4+2\pi^2\lambda^2+7\pi^4/15}{1-\iu}.
\end{align*}
\end{corollary}
\begin{proof}
Multiply \eqref{logres:eq:confluent} by $\pi t/\sin\pi t$ and
differentiate $k$ times at zero. Powers of the endpoint logarithms
are integrable, so this differentiation is legitimate. The reflection
formula and its local logarithmic expansion give
\[
 \frac{\pi t}{\sin\pi t}=\Gamma(1+t)\Gamma(1-t)
    =\exp\left(\sum_{j\ge1}\frac{\zeta(2j)}j t^{2j}\right),
 \qquad |t|<1.
\]
At $z=\iu$, the principal logarithm is
$\log(1-\iu)=\lambda$. This proves the coefficient formula and the
listed values. Only finitely many terms enter a coefficient of order
$k$, and every even zeta value is a rational multiple of $\pi^{2j}$.
\end{proof}


For a concrete weight-two pair, the real and imaginary parts of
$I_{2,2}$ give the ordinary integral identities
\begin{align}\label{logres:eq:real-second-moments}
 \int_0^1\frac{(1-u^2)\log^2(u/(1-u))}{(1+u^2)^2}\,du
 &=\frac{\log^2 2}{8}+\frac{13\pi^2}{96}
                            +\frac{\pi\log2}{8},\notag\\
 \int_0^1\frac{2u\log^2(u/(1-u))}{(1+u^2)^2}\,du
 &=\frac{\log^2 2}{8}+\frac{13\pi^2}{96}
                            -\frac{\pi\log2}{8}.
\end{align}
Both follow by expanding $\lambda^2$ in the proved formula, so their
constants are derived rather than recognized numerically.

\begin{corollary}[Two real Gaussian generating identities]
\label{logres:cor:real-generators}
For real $-1<t<1$,
\begin{align}\label{logres:eq:real-generators}
 \int_0^1\frac{(u/(1-u))^t}{1+u^2}\,du
 &=\frac{\pi\,2^{-t/2}\sin(\pi t/4)}{\sin\pi t},\notag\\
 \int_0^1\frac{u(u/(1-u))^t}{1+u^2}\,du
 &=\frac{\pi\{1-2^{-t/2}\cos(\pi t/4)\}}{\sin\pi t}.
\end{align}
At $t=0$ the removable values are respectively $\pi/4$ and
$\log2/2$. Differentiating at zero supplies every real logarithmic
moment of these two kernels.
\end{corollary}
\begin{proof}
Use \eqref{logres:eq:one} at $z=\iu$ and $z=-\iu$ and the identities
\[
 \frac1{1+u^2}=\frac12\left(\frac1{1-\iu u}+\frac1{1+\iu u}\right),
 \quad
 \frac{u}{1+u^2}=\frac1{2\iu}\left(\frac1{1-\iu u}-\frac1{1+\iu u}\right).
\]
Multiply by $\pi t/\sin\pi t$ and simplify the conjugate principal
powers. Continuity gives the zero-parameter values.
\end{proof}

The finite confluent polynomials and moment recurrence are checked exactly
in \src{verification/certify_logistic_resolvents.py}. Independent
logistic-coordinate quadratures are labeled as numerical diagnostics;
the arbitrary-order identities follow from the beta moments, partial
fractions and analytic parameter differentiation. These formulas add a
several-color integral family to the development; they do not settle the
separate mixed harmonic sums $S_6$ or $S_8$, and no literature-priority
claim is made for this classical beta-integral mechanism.


\section{Exact rational evaluation and finite certificates}
\label{ado:sec:computation}
The density zeros need not be computed in order to obtain numerical
certificates. The unsigned norm bound $\adoNorm\kappa_1\le2^{d-1}$
already gives an explicit geometric remainder for a centered
resolvent expansion. The proof works with real final order; the delivered
exact implementation accepts positive integer indices so that all
required moments are rational.

\subsection{A centered-moment expansion}
Define
\begin{equation}\label{ado:eq:centeredmoments}
 B_k=\int_0^1(2u-1)^k\kappa(u)\dd u
 =\sum_{j=0}^k(-1)^{k-j}\binom{k}{j}2^j c_{\boldsymbol{s}}(j+1).
\end{equation}
Then $|B_k|\le M$, where $M=2^{d-1}$. For $\Real z<1$, put
\[
 w=\frac{z}{2-z},\qquad |w|<1.
\]
The inequality is equivalent to $|z|<|2-z|$. The identity
$1-zu=(1-z/2)(1-w(2u-1))$ gives
\begin{theorem}[Geometric exact-evaluation certificate]
\label{ado:thm:evaluator}
For every integer $N\ge1$,
\begin{align}
 V_N(z)&=\frac{z}{1-z/2}\sum_{k=0}^{N-1}B_kw^k,
 \label{ado:eq:centeredvalue}\\
 |L(z)-V_N(z)|
 &\le\left|\frac{z}{1-z/2}\right|
          \frac{M|w|^N}{1-|w|}.
 \label{ado:eq:centerederror}
\end{align}
At $z=\iu$ and even $N=2m$, a purely rational bound is
\begin{equation}\label{ado:eq:gaussianerror}
 |L(\iu)-V_{2m}(\iu)|<\frac{2^d}{5^m}.
\end{equation}
For $|z|\le1/2$, the uniform bound is
\begin{equation}\label{ado:eq:halferror}
 |L(z)-V_N(z)|\le\frac{2^{d-1}}{3^N}.
\end{equation}
\end{theorem}
\begin{proof}
The centered geometric series converges uniformly for $u\in[0,1]$,
and its tail is bounded by $|w|^N/(1-|w|)$. Integrating against
$|\kappa|$ proves \eqref{ado:eq:centerederror}. At $z=\iu$,
$|w|=1/\sqrt5$ and $|z/(1-z/2)|=2/\sqrt5$. The multiplicative
constant is therefore $2M/(\sqrt5-1)<2M=2^d$ after extracting
$5^{-m}$. On $|z|\le1/2$, one has $|w|\le1/3$ and
$|z/(1-z/2)|\le2/3$, proving \eqref{ado:eq:halferror}.
\end{proof}
For integer indices, \eqref{ado:eq:centeredmoments} and
\eqref{ado:eq:centeredvalue} use only rational arithmetic at rational
complex $z$. An enclosing complex disk gives enclosing intervals for
each coordinate by adding and subtracting its radius. The evaluator
never treats a small residual as a proof of an identity.

\subsection{Implementation choices and replay}
The file \texttt{code/exact\_evaluator.py} first obtains the moments
by dynamic strict prefix sums. It computes $B_k$ with a common integer
denominator and a repeated-difference transform, rather than a
floating-point alternating binomial sum. Horner evaluation then gives
the exact rational center. The generic branch can upper-bound square
roots by dyadic rationals; the Gaussian and half-disk cases use the
closed rational bounds above. Near $\Real z=1$, the fixed norm-bound
precision in the generic branch can be insufficient, in which case the
implementation raises an error rather than returns an uncertified value.
It is not advertised as a complete optimal evaluator on the entire slit
plane.

From the package root, the principal replay is
\begin{verbatim}
python code/verify.py
\end{verbatim}
It reads frozen rational certificates and independently recomputes them.
The optional command
\begin{verbatim}
python code/verify.py --regenerate
\end{verbatim}
uses floating point to \emph{propose} angular brackets, then requires
rational arithmetic to certify both endpoint signs before saving them.
The proposal stage uses a strict-disk series and does not decide any
accepted sign. This discovery/replay separation is deliberate.

\subsection{Gaussian values}
The delivered Gaussian certificate file contains eight mixed-depth
examples, all evaluated with $N=240$ centered terms. Each complex-disk
radius is smaller than $10^{-81}$. The decimal entries below are
readable midpoint diagnostics; the certificates are the rational
endpoints in \texttt{data/gaussian\_enclosures.json}.
\begin{center}\small
\begin{tabular}{@{}lrr@{}}
\toprule
Index & $\Real L_{\boldsymbol{s}}(\iu)$, approximately
      & $\Imag L_{\boldsymbol{s}}(\iu)$, approximately\\
\midrule
$(2,1)$ & $-0.174715661158$ & $-0.113648917600$\\
$(1,1,1)$ & $0.099953993665$ & $-0.033577147847$\\
$(2,1,2)$ & $0.032383901801$ & $-0.026348629072$\\
$(1,2,1,1)$ & $0.001158016810$ & $0.008343470492$\\
$(2,1,1,2)$ & $0.001825088603$ & $0.006510815484$\\
$(1,1,1,1,1)$ & $-0.003396506582$ & $-0.001886779708$\\
$(2,1,2,1,1)$ & $-0.000326599679$ & $-0.000010475684$\\
$(1,1,1,1,1,1)$ & $0.000443168800$ & $-0.000335617002$\\
\bottomrule
\end{tabular}
\end{center}
In addition, the eight all-ones Gaussian values at depths one through
eight are checked against \emph{independent exact} intervals for
\[
 L_{\boldsymbol{1}_d}(\iu)=\frac{(-\tfrac12\log2+\tfrac{\iu\pi}{4})^d}{d!}.
\]
The logarithm and $\pi$ intervals are derived from different rational
series, as detailed in Section~\ref{ado:app:controls}. These eight
normalization controls are distinct from the eight stored mixed-depth
examples above.

\subsection{Certified angular brackets}
On the semicircle of radius $1/2$, set
\begin{equation}\label{ado:eq:rationalcircle}
 z(t)=\frac{1-t^2}{2(1+t^2)}+\iu\frac{t}{1+t^2},
 \qquad \theta=2\arctan t,\quad t>0.
\end{equation}
Rational $t$ gives rational real and imaginary coordinates. Every saved
root bracket has $t$-width at most $10^{-12}$; the corresponding
angle-width is less than $2\cdot10^{-12}$. With $N=104$, the
error bound \eqref{ado:eq:halferror} proves the two endpoint signs by
strict rational comparisons. For illustration:
\begin{center}\small
\begin{tabular}{@{}lcrr@{}}
\toprule
Index & $k$ & $10^{12}t_{\rm left}$ & $10^{12}t_{\rm right}$\\
\midrule
$(2,1,2)$ & 1 & 472099547153 & 472099547154\\
$(2,1,2)$ & 2 & 1487963538987 & 1487963538988\\
$(1,2,1,1)$ & 1 & 314972569528 & 314972569529\\
$(1,2,1,1)$ & 2 & 806524527583 & 806524527584\\
$(1,2,1,1)$ & 3 & 2019771345242 & 2019771345243\\
$(1,1,1,1,1,1)$ & 3 & 774596669241 & 774596669242\\
\bottomrule
\end{tabular}
\end{center}
The file \texttt{data/angular\_brackets.json} contains 24 brackets in
eight complete index families of depths two through six, together with
all exact coordinate intervals. The existence of a zero in each bracket
follows from the opposite signs. The theorem, rather than a scan,
proves uniqueness in every bracket and excludes all additional angular
zeros. For the displayed middle branch at depth six, the all-ones
formula independently gives
$\theta_3(1/2)=\arccos(1/4)$ and $t=\sqrt{3/5}$.

\subsection{What the exact suite checks}
The fresh replay reports 5,713 finite check groups:
\begin{center}\small
\begin{tabular}{@{}p{0.86\linewidth}r@{}}
\toprule
Check family & Count\\
\midrule
Strict coefficients against direct finite nested enumeration & 1080\\
Initial zero moments and leading coefficients & 120\\
Centered moments: independent binomial formula and norm bound & 2880\\
Polynomial compensation identities for arbitrary test roots & 1200\\
All-ones moments against unsigned Stirling recurrence & 240\\
Logistic-kernel derivatives through depth twelve & 11\\
Depth-three/four symbolic kernel formulas & 2\\
Exact generalized-kernel counterexample checks & 2\\
Rejected negative controls & 2\\
Rational atomic radial-argument derivative signs & 80\\
Independent exact all-ones Gaussian controls & 8\\
Replayed stored Gaussian enclosures & 8\\
Replayed angular endpoint signs & 48\\
Replayed root brackets and complete families & $24+8$\\
\midrule
Total & 5713\\
\bottomrule
\end{tabular}
\end{center}
The polynomial compensation tests deliberately use arbitrary rational
roots: they test the cancellation identity, \emph{not} positivity for
those roots. Positivity requires the actual density zeros and is proved
analytically. Similarly, finite atomic derivative checks do not replace
the interlacing-and-approximation proof of
Lemma~\ref{ado:lem:radialPick}.

\section{Independent rational logarithm and Gaussian controls}
\label{ado:app:controls}
The exact controls use only rational arithmetic and elementary remainder
bounds. They do not share the centered-moment evaluator's recurrence.
For $0<x<1$, let
\[
 A_N(x)=\sum_{j=0}^{N-1}\frac{(-1)^jx^{2j+1}}{2j+1}.
\]
The alternating-series estimate gives
\[
 \arctan x\in
 \begin{cases}
 [A_N(x),\ A_N(x)+x^{2N+1}/(2N+1)],&N\text{ even},\\
 [A_N(x)-x^{2N+1}/(2N+1),\ A_N(x)],&N\text{ odd}.
 \end{cases}
\]
The identity
\begin{equation}\label{ado:eq:machin}
 \pi=16\arctan(1/5)-4\arctan(1/239)
\end{equation}
is verified by tangent addition: if $a=\arctan(1/5)$, then
$\tan(4a)=120/119$, and
$\tan(4a-\arctan(1/239))=1$. Elementary bounds put the angle in
$(0,\pi/2)$, so it equals $\pi/4$ and not another tangent branch.
The implementation uses $N=120$ and $N=40$ for the two arctangent
series.

For the logarithm, set
\[
 B_N=2\sum_{j=0}^{N-1}\frac{(1/3)^{2j+1}}{2j+1}.
\]
Integrating the geometric series for $2/(1-x^2)$, or using the
power series of $2\operatorname{arctanh}x$, gives
\begin{equation}\label{ado:eq:log2interval}
 B_N<\log2<
 B_N+\frac{2(1/3)^{2N+1}}{(2N+1)(1-1/9)}.
\end{equation}
The code uses $N=160$. Both interval endpoints are exact rational
numbers. Cartesian complex interval multiplication then encloses
$(-\tfrac12\log2+\tfrac{\iu\pi}{4})^d/d!$.

For each depth $1\le d\le8$, the independently generated real and
imaginary intervals lie wholly inside the corresponding coordinate
intervals returned by the centered evaluator with 240 terms. This is
an exact containment check, not a decimal comparison. It is a strong
normalization and branch control, while the identity
\eqref{ado:eq:allonesLi} still rests on its differential-equation proof.

The exact evaluator also supplies a dyadic upper bound on $\sqrt{x}$
for a nonnegative rational $x=a/b$. With an integer precision $p$, set
\[
 q=\frac{\lfloor\sqrt{\lfloor 2^{2p}a/b\rfloor}\rfloor+1}{2^p}.
\]
Then $q^2>x$. This is computed by integer square root and exact
comparison. If its use in the generic geometric-tail formula fails to
produce an upper bound smaller than one for $|w|$, the code stops
rather than accepting a possibly divergent bound.

The complete original source and delivered rational certificates remain
under \src{reports/angular-zero-continuation/continuations/all-depth-signed-oscillation/}.
The isolated replay record in \src{verification/sixth-replay/all-depth/}
checks the finite arithmetic independently of the analytic zero-count proof.
The manuscript uses the stronger positive-transform radial theorem proved
next, while retaining the incoming interlacing argument as an independent
route. Fractional leading indices at higher depth still require a new
construction: iterating an integer number of moment operators does not
establish a fractional variation theorem.
